\section{First resonance and the correct finite part}
\label{dhr:sec:resonance}

Fix $r\geq1$, initially $a,c>0$, and set
\begin{equation}\label{dhr:eq:resonant-parameters}
 b_0=a+c+2r-2,\qquad b=b_0+u.
\end{equation}
Let $L(u)$ be any holomorphic function near zero with $L(0)=1$.
It is a normalization of \emph{each} particle's weight. Define
\begin{equation}\label{dhr:eq:amplitude-E}
 \begin{split}
 A(u)&=L(u)^r Z_{r-1}(a,c;b_0+u),\\
 E(u)&=\frac{\Gamma(r)\Gamma(a+r-1)\Gamma(c+r-1)\Gamma(1+u)^2}
 {\Gamma(r+u)\Gamma(a+r-1+u)\Gamma(c+r-1+u)}.
 \end{split}
\end{equation}
Both are holomorphic near zero, and $E(0)=1$. Although $A$ uses the
$(r-1)$-particle partition, its normalization is $L^r$, not $L^{r-1}$:
the escaping largest index contributes a normalization factor as well.

Write the contribution of a prescribed largest index as
\begin{equation}\label{dhr:eq:maximum-mass}
 \begin{split}
 K_n(u)={}&L(u)^r w_n(a,c;b_0+u)\\
 &\quad\times\sum_{0\leq m_1<\cdots<m_{r-1}<n}
 \Delta(\bm m)^2\prod_{i=1}^{r-1}
       (n-m_i)^2w_{m_i}(a,c;b_0+u).
 \end{split}
\end{equation}
For $r=1$ the inner sum is one, including at $n=0$; for $r>1$ it
vanishes if $n<r-1$. Each $K_n$ is holomorphic near zero.

\begin{theorem}[Largest-index subtraction]\label{dhr:thm:resonance}
For every fixed $\alpha$ with $\Rea\alpha>0$, the series
\begin{equation}\label{dhr:eq:D-def}
 D(u;\alpha)=\sum_{n\geq0}
 \left\{K_n(u)-A(u)(n+\alpha)^{-1-u}\right\}
\end{equation}
is locally normally convergent near $u=0$, and
\begin{equation}\label{dhr:eq:D-master}
 \boxed{D(u;\alpha)=A(u)
       \left\{\frac{E(u)}u-\zeta(1+u,\alpha)\right\}.}
\end{equation}
The apparent pole on the right is removable. Every fixed derivative
in $u$ can be taken termwise. The assertions hold jointly locally
holomorphically for $\Rea a,\Rea c,\Rea\alpha>0$.
\end{theorem}

\begin{proof}
Dividing the products in Theorem~\ref{dhr:thm:gauss} at particle
numbers $r$ and $r-1$, with the same $b=b_0+u$, gives
\begin{equation}\label{dhr:eq:meromorphic-master}
 L(u)^r Z_r(a,c;b_0+u)=\frac{A(u)E(u)}u.
\end{equation}
For instance, the ratio of the factors common to the two products
contains
\[
 \prod_{j=0}^{r-2}
 \frac{\Gamma(r-1+u-j)}{\Gamma(r+u-j)}
 =\frac{\Gamma(1+u)}{\Gamma(r+u)}.
\]
The remaining factor has $\Gamma(u)=\Gamma(1+u)/u$. These two
factors explain the square of $\Gamma(1+u)$ in $E$.
Equation~\eqref{dhr:eq:meromorphic-master} initially holds for
$\Rea u>0$ sufficiently small.

We next justify the subtraction independently of that product's
meromorphic continuation. Work on a compact parameter set and a small
closed disk $|u|\leq\delta<1/2$. The absolute $(r-1)$-particle sum
has a uniformly finite first moment:
\begin{equation}\label{dhr:eq:lower-first-moment}
 \sum_{m_1<\cdots<m_{r-1}}
 (1+m_1+\cdots+m_{r-1})
 |\Delta(\bm m)|^2\prod_i|w_{m_i}|<C.
\end{equation}
Indeed, the largest exponent in the separate-index majorant is
$-3-\Rea u$ before inserting a first moment, and at most
$-2+\delta<-1$ after insertion. This also proves the assertion
uniformly for complex $a,c$ in compact right-half-plane sets.

For $0\leq m_i<n$,
\[
 \left|1-\prod_i(1-m_i/n)^2\right|
 \leq\frac{2}{n}\sum_i m_i.
\]
The missing region $m_{r-1}\geq n$ is likewise bounded by the first
moment divided by $n$. It follows that the inner sum in
\eqref{dhr:eq:maximum-mass}, divided by $n^{2r-2}$, is
\[
 Z_{r-1}(a,c;b_0+u)+O(n^{-1})
\]
uniformly on that compact set. The assertion is immediate for $r=1$.
The uniform gamma-ratio expansion gives
\[
 w_n(a,c;b_0+u)=n^{-2r+1-u}\{1+O(n^{-1})\}.
\]
Consequently, for $n\geq1$,
\begin{equation}\label{dhr:eq:asymptotic-mass}
 K_n(u)=A(u)n^{-1-u}+O(n^{-2+\delta}).
\end{equation}
Replacing $n$ by $n+\alpha$ in the leading term changes only the
summable remainder, uniformly on compact $\alpha$-sets. The finitely
many initial terms are harmless. This proves local normal convergence
of \eqref{dhr:eq:D-def}. Cauchy's integral formula on a slightly
larger parameter disk proves normal convergence for every fixed
$u$-derivative; there is no unproved exchange of an infinite series
with differentiation at a pole.

For $\Rea u>0$, sum the leading counterterm by Hurwitz zeta and use
\eqref{dhr:eq:meromorphic-master}. This yields
\eqref{dhr:eq:D-master} on a nonempty open set. Both sides extend
holomorphically through zero, because $E(0)=1$ and the Hurwitz residue
is one. The identity theorem completes the proof, including joint
holomorphy in the auxiliary parameters.
\end{proof}

Put $A_j=A^{(j)}(0)$ and
\begin{equation}\label{dhr:eq:chi}
 \chi=E'(0)=2\psi(1)-\psi(r)-\psi(a+r-1)-\psi(c+r-1).
\end{equation}
In particular,
\begin{equation}\label{dhr:eq:residue}
 A_0=\prod_{j=0}^{r-2}
 \frac{j!\,\Gamma(a+j)\Gamma(c+j)\Gamma(r-j)}
 {\Gamma(c+2r-2-j)\Gamma(a+2r-2-j)}.
\end{equation}

\begin{corollary}[Hard-cutoff constant]\label{dhr:cor:cutoff}
For $a,c>0$,
\begin{equation}\label{dhr:eq:cutoff-constant}
 \boxed{\lim_{N\to\infty}\left[
 \sum_{0\leq n_1<\cdots<n_r\leq N}
   \Delta(\bm n)^2\prod_i w_{n_i}(a,c;b_0)
       -A_0\log N\right]=A_0\chi.}
\end{equation}
The corresponding spectral finite part is
\begin{equation}\label{dhr:eq:spectral-vs-cutoff}
 \FP_{u=0}\{L(u)^rZ_r(a,c;b_0+u)\}=A_1+A_0\chi.
\end{equation}
Thus the hard-cutoff constant is obtained from the spectral finite
part by subtracting $A_1$.
\end{corollary}
\begin{proof}
At zero, Theorem~\ref{dhr:thm:resonance} gives
\[
 D(0;\alpha)=A_0\{\chi-\gamma_0(\alpha)\}
            =A_0\{\chi+\psi(\alpha)\}.
\]
The partial sum of the subtracted counterterm is
$A_0\{\psi(N+1+\alpha)-\psi(\alpha)\}$.
Since $\psi(N+1+\alpha)=\log N+o(1)$, the auxiliary $\alpha$
cancels and proves \eqref{dhr:eq:cutoff-constant}.
Expanding $A(u)E(u)/u$ directly gives
\eqref{dhr:eq:spectral-vs-cutoff}.
\end{proof}

\begin{remark}[Why the normalization matters]
Changing $L$ while preserving $L(0)=1$ leaves the entire cutoff sum
unchanged. It changes $A_1$ by $rL'(0)A_0$ relative to the choice
$L\equiv1$, and changes the spectral finite part by the same amount.
Therefore identifying spectral and cutoff constants without specifying
the amplitude derivative cannot be a normalization-invariant statement.
This is a correction to that possible extrapolation, not an allegation
that the inspected manuscript makes the identification.
\end{remark}
